\cotefile{157-4}
\scanpage{50}
\mspage{43}

Ces conditions \unsure{indiquent|impliquent} que pour tout $I$ dans \uline{$\mathrm{Cat}$}, il existe \unsure{deux|une} str. de Quillen \unsure{faible|faibles} \unsure{exacte|exactes} \unsure{à|:} g et \unsure{close} \unsure{pour} $W(I)$ dans $M(I)$. \deleted{\unsure{donnée}} Une telle structure se \unsure{définit|définirait} en termes d'une structure similaire $C,\mathrm{TC}=C\cap W,F=(\unsure{\mathrm{TC}|\mathrm{TT}})_{\unsure{*|\alpha|d}},$ $\mathrm{TF}=C_*)$ pour $W$, par

\[ C(I)=\{f\in\mathrm{Fl}(M(I))\mid\alpha_s^{\unsure{*}}(f)\in C\quad\forall s\in\mathrm{Ob}\,I\}, \]

puis

\[ \begin{aligned}
\mathrm{TC}(I)&=C(I)\cap W(I)=\{f\in\mathrm{Fl}(M(I))\mid\alpha_s^*(f)\in\mathrm{TC}\ \forall s\in\mathrm{Ob}\,I\},\\
F(I)&=(\mathrm{TC}(I))_*\\
\mathrm{TF}(I)&=C(I)_*\quad .
\end{aligned} \]

\uline{Corollaire}\inserted{\unsure{- scholie}} Ces conditions impliquent que pour tout \deleted{$\unsure{W(I),)}$} $I$ dans $\mathrm{Cat}$, $(M(I),W(I))$ donne lieu à une \inserted{bonne} théorie des fibres et cofibres $W(I)$-homotopiques, avec carrés $W(I)$-cartésiens \deleted{\unsure{et}} $W(I)$-cocart., et carrés $W(I)$-cocartésiens $W(I)$-cart. -- donnant des suites \uline{exactes} longues des \unsure{cosuspensions} et des suspensions, dans le cas où $M$, donc aussi $M(I)$, est une catégorie avec objet \unsure{neutre}.
% archive sheet number: 49
\scanpage{51}
\mspage{44}

Je m'occuperai de la démonstration des théorèmes précédents, ainsi que du th. de factorisation du § 1, quand j'aurai mes notes sur les catégories \deleted{\unsure{strictement|stratifiées}} accessibles et filtrations \unsure{ordinales|cardinales}. En attendant, j'ai envie déjà d'exploiter un peu plus le formalisme auquel on parvient moyennant les hypothèses faites ci-dessus. (Je suis à peu près sûr tout au moins que la condition \unsure{a$'$)|a$''$)} suffit à impliquer \deleted{toutes} les autres. Et c'est \deleted{\unsure{bien}} sous la forme \unsure{a$''$)|a$'''$)} que \deleted{\unsure{cette}} conditions vont se vérifier en pratique, je présume, notamment dans le cas où $M\unsure{:=|=}\mathrm{Cat}$.

4. \uline{Les \unsure{variances}} $u^*,u_*,u_!$ \uline{pour} $(M(I),W(I))$ ($I$ variable).

Je suppose \deleted{les \unsure{conditio}} les conditions du § précédent vérifiées pour $(M,W)$\inserted{(cf th. p. 42)}, et choisi \unsure{une} structure de $\unsure{\mathrm{Qf}|\mathrm{QF}}$ close et exacte à gauche $(C,\mathrm{TC},F,\mathrm{TF})$ pour $W$, d'où pour \unsure{tt} $I$ dans $\mathrm{Cat}$
% archive sheet number: 50
\scanpage{52}
\mspage{45}

une structure $\unsure{\mathrm{Qf}|\mathrm{QF}}$ close et exacte à gauche $(C(I),\mathrm{TC}(I),F(I),\mathrm{TF}(I))$ pour $W(I)$, dans $M(I)$, décrite page 43. Soit \deleted{\unsure{d'autre part}} \inserted{\unsure{maintenant}}

\[ u:I'\longrightarrow I \]

une flèche dans $\mathrm{Cat}$, d'où une flèche

\[ u^*:(M(I),W(I))\longrightarrow(M(I'),W(I')). \]

On sait, vu que $M$ est stable par $\varinjlim$ \unsure{quelc.} et aussi (étant de plus accessible) par $\varprojlim$ quelconques, que $u^*$ admet des foncteurs adjoints à gauche et à droite

\[ u_!,u_*:\quad M(I')\longrightarrow M(I), \]

et je me propose ici de \deleted{\unsure{donner}}\inserted{résumer} \deleted{\unsure{les}} propriétés les plus immédiates de $u^*$ et de $u_!,u_*$, p.r. aux structures de $\mathrm{Qf}$ envisagées.

\uline{Proposition 1} : Avec les notations précédentes, on a les relations

\[ u^*(C(I))\subset C(I'),\quad u^*(\mathrm{TC}(I))\subset\mathrm{TC}(I') \tag*{(1)} \]

\[ u_*(F(I'))\subset F(I),\quad u_*(\mathrm{TF}(I')\subset\mathrm{TF}(I) \tag*{(2)} \]

\[ \begin{aligned}
(u^*(F(I)\unsure{)|]}\subset F(I'))&\Longleftrightarrow u_!(\mathrm{TC}(I'))\subset\mathrm{TC}(I)\\
(u^*(\mathrm{TF}(I))\subset\mathrm{TF}(I'))&\Longleftrightarrow u_!(C(I'))\subset C(I)
\end{aligned} \tag*{(3)} \]

D'autre part, \unsure{je} ne \deleted{\unsure{puis}} sais rien dire de positif pour les inclusions hypothétiques

\[ \begin{aligned}
u_{\unsure{*}}(C(I'))\subset C(I),&\quad u_*(\mathrm{TC}(I'))\subset\mathrm{TC}(I)\\
u_!(F(I'))\subset F(I),&\quad u_!(\mathrm{TF}(I'))\subset\mathrm{TF}(I),
\end{aligned} \]

si ce n'est que, pour un $u$ tant soit peu général (\unsure{et|ne} pas un iso, disons), $u_*$ n'a aucune envie
% archive sheet number: 51
\scanpage{53}
\mspage{46}

d'être compatible avec les structures $C,\mathrm{TC}$ (propriété d'exactitude \inserted{homotopique} à droite, \unsure{moralement|normalement}, alors que $u_*$ est exact à gauche et \deleted{\unsure{p}} en général pas du \deleted{tout} à droite), ni $u_!$ avec les structures $F,\mathrm{TF}$ (commentaire dual à donner, $u_!$ est exact à \uline{droite}, et n'a pas envie d'être ``homotopiquement exact à gauche'').

Démonstration de \unsure{(1) (2) (3)}. Les relations (1) sont évidentes par définition. Les relations \unsure{(1)|(2)}, (3) le sont à peine moins, vu la propriété \unsure{surtout|somme toute} générale suivante, dont la place serait au § 1 :

\uline{Proposition 2} Soit\deleted{\unsure{ent}}

\diagram{\begin{tikzcd}[ampersand replacement=\&]
C' \arrow[r,shift left,no head,"u_!"] \& \deleted{\unsure{C}}C\arrow[l,shift left,"u^*"]
\end{tikzcd}}

un couple de foncteurs adjoints $(u_!,u^*)$, et soit $\Phi'\subset\mathrm{Fl}\,C'$, $\Psi\subset\mathrm{Fl}(C)$. Alors

\[ (u_!(\Phi')\longleftrightarrow\Psi)\Longleftrightarrow(\Phi'\longleftrightarrow u^*(\Psi)) \]

\uline{Corollaire} (C'est le $=$ énoncé, avec seulement \unsure{des|aux|un} \unsure{notation|notations} changée\unsure{s}). Supposons que $u^*:C\to C'$ \unsure{admette} un adjoint à \uline{droite} $u_*$ (\unsure{pas} forcément un adjoint à gauche).

\diagram{\begin{tikzcd}[ampersand replacement=\&]
C'\arrow[r,shift right,no head,"u_*"'] \& C\arrow[l,shift right,"u^*"']{,}
\end{tikzcd}}

soit $\Phi\subset\mathrm{Fl}(C)$, $\Psi'\subset\mathrm{Fl}(C')$, alors

\[ (u^*(\Phi)\longleftrightarrow\Psi')\Longleftrightarrow(\Phi\longleftrightarrow u_*(\Psi')) \]

C'est la prop. 2 \deleted{appliquée \unsure{au}} \inserted{\unsure{via|vu|en}} \deleted{$\unsure{C},C^{\unsure{n}}$}

\diagram{\begin{tikzcd}[ampersand replacement=\&]
\deleted{C'_0}\arrow[r,shift left,"\deleted{u_{0!}}"] \& \deleted{C_0}\arrow[l,shift left,"\deleted{u_0^*}"]
\end{tikzcd}}

\deleted{\unsure{tous les objets}} avec $C'_0=C$, $C_0=C'$, $(u_{0!}:C\to C')\deleted{\unsure{=u}}=u^*$, et $(u_0^*:C'\unsure{-|\to}C)=u_*$.
% archive sheet number: 52
\scanpage{54}
\mspage{47}

Pour les \deleted{équivalences (2)} inclusions (2), on utilise le corollaire \inserted{de la prop. 2} pour les couples $\Phi=\mathrm{TC}(I),\Psi'=F(I')$, resp. $\Phi=C(I)$; $\Psi'=\unsure{E|F|\mathrm{TF}}(I')$, et on applique (1), qui donne $u^*(\Phi)\leftarrow\Psi'$, d'où $\Phi\longleftrightarrow u_*(\Psi')$, ce qui n'est autre que (2). Pour les équivalences (3), on utilise la proposition 2, en faisant $\Psi=F(I),\Phi'=\mathrm{TC}(I')$ resp\unsure{t}. $\Psi=\mathrm{TF}(I),\Phi'=C(I')$.

\uline{Corollaire 1} (de la proposition 1) : Supposons $u:I'\unsure{-|\to}I$ telle que pour tout $i\in\mathrm{Ob}\,I$, la catégorie $i\backslash I'$ soit discrète\textsuperscript{*} \deleted{(p. ex. \unsure{que $u$ soit} \unsure{isom.} local)} (p. ex. $I'$ discrète). Alors on a

% slanted

\margin{Plus généralement, il suffit que que \unsure{la|les|toute} \unsure{composante|composantes} \unsure{connexe|connexes} de $i\backslash I'$ ait \inserted{$e_\alpha$} un objet initial, de sorte que \[ \underset{i\backslash I'}{\varinjlim}\,X'(i')\simeq\coprod_{\alpha\in\pi_0(i\backslash I')}X'(e_\alpha). \] C'est le cas notamment si $I'\to I$ est un isom. local, p. ex. une \unsure{immers. ouverte}.}

\[ \left\{\begin{aligned}
u_!\deleted{\unsure{(}}C(I')&\subset\deleted{\unsure{T}}C(I)\unsure{)},&u_!(\mathrm{TC}(I'))&\subset\mathrm{TC}(I)\\
u^*(F(I))&\subset F(I'),&u^*(\mathrm{TF}(I))&\subset\mathrm{TF}(I')
\end{aligned}\right. \]

En effet, les relations de la deuxième ligne équivalent (à cause de (3)) à celles de la première. Il faut prouver que si $f$ dans $\mathrm{Fl}\,M(I')$ est dans $C(I')$ resp. $\mathrm{TC}(I')$, $u_!(f)$ est dans $C(I)$ resp. $\mathrm{TC}(\unsure{I|II})$, ce qui signifie aussi que $\alpha_s^*(u_!(f))$ est dans $C$ resp. $\mathrm{TC}$, pour $\forall\,s\in\mathrm{Ob}\,I$. Or le foncteur

\[ \alpha_s^*u_!:\deleted{\unsure{A}}M(I')\longrightarrow M \]

s'explicite par la formule bien connue

\[ \alpha_s^*(u_!(X'))\overset{\text{déf}}{=}u_!(X')(s)\simeq\varinjlim_{i'\in s\backslash I'}X'(i').\quad\text{Comme} \]
% archive sheet number: 53
\scanpage{55}
\mspage{48}

\deleted{\unsure{on}} $s\backslash I'$ est discrète, on a donc l'isom. fonctoriel en $X'$ dans $M(I')$

\[ \alpha_s^*(u_!(X'))\simeq\coprod_{\substack{(s'\unsure{,|;}\unsure{h|\lambda}:s\to u(s'))\\\text{dans }\mathrm{Ob}\,s\backslash I'}}X(s'), \]

\deleted{\unsure{donc}} donc \inserted{si $f:X'\unsure{-|\to}Y'$,} $\alpha_s^*u_!(f)$ s'explicite comme une somme directe \inserted{dans $\mathrm{Fl}(M)$} de flèches de la forme

\[ f_{s'}:X'(s')\longrightarrow Y'(s'), \]

lesquelles sont dans $C$ resp. $\mathrm{TC}$ puisque $f$ est dans $C(I')$ resp. $\mathrm{TC}(I')$. Comme $C,\mathrm{TC}$ sont stables par sommes directes quelconques. ($C,\mathrm{TC}$ étant clos à gauche), on gagne.

\uline{Corollaire 2} Soit $I$ dans $\mathrm{Cat}$, \deleted{Alors \unsure{pour $\forall\,s$}} $f\in\mathrm{Fl}(M(I))$. Pour que $f$ soit dans $F(I)$ (resp. \deleted{dans} $\mathrm{TF}(I)$), il \uline{faut} que $\forall s\in\mathrm{Ob}\,I$, $f_s:X(s)-Y(s)$ le soit. (Mais cette condition n'a pas \unsure{du|un} \deleted{\unsure{tout}} l'air d'être suffisante !) En particulier, soit $X\in\mathrm{Ob}\,M(I)$, alors $X\in\mathrm{Ob}\,M(I)_f$ $\Longrightarrow X(s)\in\mathrm{ob}\,M_f\quad\forall s\in\mathrm{Ob}\,I$.

On applique le cor. 1 à $\alpha_{\unsure{e|s}}:e-I$.

\uline{Scholie} La prop. 1 \inserted{(\unsure{relations|relation} (1))} montre que le foncteur $u^*$ transforme carrés \inserted{$W$-}cocartésiens en \unsure{itou|idem}, suite exacte longue de suspensions en \unsure{itou|idem} -- c'est ``l'exactitude à droite ``homotopique'' de $u^*$. On s'attend à ce qu'il en soit de $=$ pour les carrés cartésiens et les suites exactes
% archive sheet number: 54
\scanpage{56}
\mspage{49}

longues de \unsure{cosuspension|cosuspensions}, mais pour le prouver il faudra un travail, qu'on fera dans la suite.

Les relations (2) montrent que le foncteur $u_*$ transforme carrés $W$-cartésiens en \unsure{itou|idem}, donc suites exactes longues de cosuspensions en \unsure{itou|idem}. Je ne m'attends \uline{pas} à ce qu'il en soit de $=$ \inserted{en général} pour les carrés $W$-cocartésiens, -- i.e. que $u_*$ soit ``homotopiquement exact à droite'', et non seulement à gauche. Mais il faudrait au moins que je fasse des exemples (avec $I',I$ des cat finies ordonnées).

Quant à $u_!$, la proposition ne dit rien de bien palpable à son sujet. Néanmoins on s'attend à ce que $u_!$ transforme carrés $W$-cocartésiens en \unsure{itou|idem} -- c'est la moindre des choses ! Ou peut-être \deleted{je} fais erreur, en confondant les foncteurs $u_*,u_!$ envisagés ici, avec les foncteurs similaires \deleted{\unsure{sur}}

\diagram{\begin{tikzcd}[ampersand replacement=\&]
\mathrm{HO}_W(I')\arrow[r,shift left,"u_!^W"]\arrow[r,shift right,"u_*^W"'] \& \mathrm{HO}_W(I){,}
\end{tikzcd}}

adjoints à g. resp. à dr. de

\diagram{\begin{tikzcd}[ampersand replacement=\&]
\mathrm{HO}_W(I') \& \mathrm{HO}_W(I)\arrow[l,"u_W^*"']
\end{tikzcd}}

induit par

\[ u^*:(M(I),W(I))\longrightarrow(M(I'),W(I')). \]

Mais en général $u_*,u_!$ ne sont pas compatibles
% archive sheet number: 55
\scanpage{57}
\mspage{50}

avec les localiseurs $W(I')$, $W(I)$, donc ils ne \deleted{peuvent} doivent pas être confondus avec les foncteurs $u_{\unsure{!|*}}^W,u_*^W$. \deleted{\unsure{précédent}} Donc j'ai été peut-être hâtif, dans ce que j'ai dit au sujet de $u_*$ et de $u_!$. Je serai obligé d'étudier la situation de très près, quand j'essayerai de construire les foncteurs $u_*^W$ et $u_!^W$. Pour le moment, je n'ai encore aucune assurance que les conditions envisagées \unsure{sur|pour} $(M,W)$ \inserted{suffisent à} \unsure{assurer|assurent} l'existence des $u_*^W,u_!^W$ -- pas $=$ pour $I,I'$ des ens. ordonnés finis !

% margin note written sideways

\margin{Comparer avec le lemme p. \unsure{63|13|15} et sa démonstration (plus exactement, le cor. du lemme).}

% slanted

\margin{\uline{Prop. 3} Soit $X_{\unsure{\bullet|.}}$ dans $M(\Delta^1)$, $X=(X_0\xleftarrow{\alpha}X_1)$. Alors $X_{\bullet}$ est dans \deleted{\unsure{$\mathrm{ob}.$}} $M(\Delta^1)_f$ \unsure{sss|ssi} $X_0,X_1$ dans $M_f$ et $\alpha\in F$.}

\margin{\deleted{\uline{\unsure{Proposition 3}}}\inserted{\unsure{Cor. à la|Cor. à 4|Cor. 1}} Soit $X\in M(I)_f$, alors pour \unsure{tte} $\alpha:s\to t$, dans $I$ qui est \uline{mono}, $X(\alpha):X(\unsure{t|f})\to X(s)$ est dans \uline{$F$}, (On applique cor. \uline{1} \inserted{(p. 47)} à $\widetilde{\alpha}:\Delta^1\to I$)}

\margin{[\uline{Cor.} ? Soit $\unsure{I'|II'}\overset{u}{\leftrightarrow}I$ \uline{imm. ouverte}. Alors \unsure{pour} $X\in\mathrm{ob}\,M(I)_f$, $\varprojlim_I X(i)\to\varprojlim_{I'}\unsure{\alpha|X}(i')$ est dans $F$.}

\textcircled{5}. \uline{Catégories de modèles quillénisables}.

Je reviens au cas général d'une catégorie de modèles $(M,W)$.

\llap{\raisebox{-1.75\baselineskip}[0pt][0pt]{\rule{0.4pt}{1.85\baselineskip}}\hspace{0.5em}}\uline{Définition 1} On dit que $(M,W)$ est \uline{quillénisable} s'il existe une structure de $\unsure{\mathrm{Qf}|\mathrm{QF}}$ \uline{exacte} $(C,\mathrm{TC},F,\mathrm{TF})$ pour $M$.

C'est là une condition autoduale, mais que je ne sais vérifier en pratique que moyennant l'exactitude soit à gauche, soit à droite -- donc par une condition plus forte qui détruit la symétrie. C'est pour la sauvegarder dans la mesure du possible que j'introduis la définition précédente, ainsi que la

\uline{Définition 2} Soit, en plus de $(M,W)$, $\deleted{\mathrm{Diag}}\subset\mathrm{Cat}$ une sous-catégorie pleine de $\mathrm{Cat}$. On dit que
% archive sheet number: 56
\scanpage{58}
\mspage{51}

$W$ est \uline{quillénisable sur $\mathrm{Diag}$}, si pour tout $I$ dans $\mathrm{Diag}$, \deleted{\unsure{le}} $(M(I),W(I))$ est quillénisable.

% slanted

\margin{Je supposerai que $\mathrm{Diag}$ contient \deleted{\unsure{tout}} au moins les cat. provenant d'ens\unsure{.} ordonnés finis.}

En pratique, j'utiliserai cette \unsure{locution|notion} dans le cas où $\mathrm{Diag}=\mathrm{Cat}$, ou alors, qu'il \unsure{satisfasse} au moins les propriétés de stabilité utiles pour le domaine d'un dérivateur : stable par \deleted{\unsure{prod}}\,$\varprojlim$ finies, par sommes finies -- éventuellement par $\underset{\unsure{-|\to}}{\lim}$ finies, etc.

\uline{Exemple \uline{1}} : Celui des § 3, 4, où on a dégagé des conditions sur $(M,W)$, avec $M$ une ``\unsure{grosse|pure}'' catégorie et $M$ \unsure{accessible}, pour que \unsure{lui|$W$} soit quillénisable sur $\mathrm{Cat}$ tout \unsure{entier|entière}. C'est le cas qui me paraît à présent le plus intéressant.

\uline{Exemple 2} : Soit\deleted{\unsure{ent}} $A$ une catégorie abélienne, et $M=K^{\unsure{\bullet|\cdot|\prime}}(A)$ la catégorie des complexes de \unsure{cochaines|cochaînes} à valeurs dans $A$, $W$\inserted{$=W_A$} l'ens. des quasi-isomorphismes, donc

\[ MW^{-1}=K^\bullet(A)\,W_A^{-1}\simeq D(A)\quad\text{catégorie dérivée de }A. \]

On définit \unsure{aussi} $M^+,M^-,M^b$, avec les localiseurs induits $W_A^+,W_A^-,W_A^b$, et les cat. de fractions $D^+(A),D^-(A),D^b(A)$. On
% archive sheet number: 57
\scanpage{59}
\mspage{52}

aura, pour $I$ dans $\mathrm{Cat}$

\[ \begin{aligned}
M(I)&=K^{\unsure{\bullet|\cdot}}(A(I))\\
W(I)&=\underbrace{W_{A(I)}}_{\text{quasi-isomorphismes}}\unsure{\subset|\in}\mathrm{Fl}(K^\bullet(A(I)))
\end{aligned} \]

et de même pour $M^\alpha(I),W^\alpha(I)$ $(\alpha\in\{+,-,b\})$, du moins \uline{si $I$ est fini}. Il y aussi les sous-catégories $\unsure{K^{\geq0}|K^{>0}}(A)$ et $K^{\leq0}(A)$ des complexes de \unsure{cochaines|cochaînes} resp. de \unsure{chaines|chaînes} au sens strict (avec $K^i=0$ si $i\geq0$, resp\unsure{t}. si $i\leq0$). On aura pour $I$ quelconque

\[ \left\{\begin{aligned}
K^{\geq0}(A)(I)&\simeq K^{\geq0}(A(I))\\
W_A^{\geq0}(I)&=W_{A(I)}^{\geq0}
\end{aligned}\right. \]

et de $=$ pour $K^{\leq0},W^{\leq0}$.

D'après Quillen \unsure{pag.|p.} 1.2, $(K^+(A),W_A^+)$ est quillénisable si $A$ a assez d'injectifs, $(K^-(A),W_A^-)$ l'est si $A$ a assez de projectifs. Il \unsure{ne dit|n'en est} rien pour $(K^\bullet(A),W_A^{\unsure{*|\bullet}})$ sans plus -- pourtant on fait de l'algèbre cohomologique avec \deleted{$K(A),W_A)$} $D^{\unsure{?|\bullet}}(A)$ sans plus, et $=$ sans faire d'hypothèse d'existence d'injectifs ou de projectifs. J'ai un peu oublié ce que Verdier a fait à ce sujet -- je devrais quand même regarder sa thèse (!), ou bien l'exposé de Hartshorne,
% archive sheet number: 58
\scanpage{60}
\mspage{53}

Toujours est-il qu'on a là des exemples de type \inserted{tout} différent de ceux envisagés dans 1°. Ainsi, il n'y a pas à supposer que $M$ \deleted{soit} stable par \inserted{petites} $\varinjlim$ -- $M$ peut même être une petite catégorie, \inserted{elle l'est} \deleted{\unsure{si}} $A$ elle-même l'est. Si $A$ est stable par \deleted{\unsure{petites}}\inserted{(petites)} $\varinjlim$ filtrantes, et est accessible (p. ex. a une petite famille génératrice, et les $\varinjlim$ filtrantes sont exactes) alors $K^{\unsure{\bullet|\cdot}}(A)$ \uline{etc} ont les $\unsure{=\mathrm{s}|=|\Rightarrow}$ propriétés, et $(K^+(A),W_A^+)$ est quillénisable sur $\mathrm{Cat}$ \deleted{tout} entier. (C'est le cas où on sait que $A$ a assez d'injectifs, et l'algèbre \unsure{homologique|cohomologique} se fait tout à coups de résolutions injectives.)

% slanted

\margin{!! attention, $K^+(A)$ n'est \uline{pas} stable par $\varinjlim$ filtrantes infinies -- \unsure{d'où|et donc} \unsure{ne|on} peut appliquer les acquis des § 2, 3 \unsure{directement|seulement} à $K^{\unsure{+|\geq0}}$}

Dans ce cas on sait que les \deleted{\unsure{$u$}}

\[ u_*:D^+(I',A)\longrightarrow D^+(I,A) \]

existent (comme adjoints à droite de

\[ u^{\unsure{*}}:D^+(I,A)\longrightarrow D^+(I',A)), \]

mais je ne suis plus trop sûr si on sait définir $u_*$ aussi sur $D(I',A)$ \deleted{tout} entier, et je doute que $u_!$ existe sans hypothèse sur $u$ (genre : dimension coh. finie). Sans doute, on doit tjrs avoir

\[ u_!:D^-(I',A)\longrightarrow D^-(I,A), \]
% archive sheet number: 59
\scanpage{61}
\mspage{54}

adjoint à gauche de

\[ u^*:D^-(I,A)\longrightarrow D^-(I'',A). \]

Mais je \deleted{\unsure{ne}} ne sais plus trop s'il faut s'attendre à pouvoir définir $u_!$ sur $D(I',A)$ \deleted{tout} entier, \deleted{\unsure{au}} \deleted{\unsure{en}} dehors des hypothèses genre dim. coh. finie. La triste vérité, c'est que j'ai $\pm$ oublié la théorie des foncteurs dérivés ! (Il faudrait que je demande \unsure{encore} à Jean Malgoire de m'envoyer des textes canoniques.)

\uline{Scholie} Si $(M,W)$ est quillénisable sur \deleted{\unsure{$\mathrm{Cat}$}} $\mathrm{Diag}\subset\mathrm{Cat}$, \inserted{et $=$ sans qu'il soit quillénisable,} il définit un prédérivateur \deleted{sur} $\mathbb{D}_{(M,W)}$ sur $\mathrm{Diag}$, \deleted{\unsure{qui est}} On \unsure{posera|écrira} parfois

\[ \mathrm{HO}_{M,W}(I)\ \text{ou}\ \mathrm{HO}_W(I)=\mathbb{D}_{M,W}(I). \]

L'hypothèse quillénisable implique, sur \inserted{\unsure{du moins pour|de même que}} ces ``catégories homotopiques, la structure \unsure{(relativement)} \unsure{``faible''} ``triangulée'', provenant des suites ``\uline{exactes}'' de suspension et de cosuspension\textsuperscript{*}. (La condition Quillénisable \unsure{assure} l'exactitude de ces suites \ldots)

% slanted

\margin{dans le cas \unsure{(particulier)|(cf. chapitre)} où $M$ est \unsure{pointée}, \unsure{sinon}, \unsure{par hypoth.|pour objet} \unsure{d'additivité} \illegible{1}, il y a \unsure{la structure|les structures} \unsure{associée|associées} aux \deleted{\unsure{foncteurs|carrés}} carrés cartésiens \unsure{resp.|et} cocartésiens. \unsure{Il suffit d'|Il s'agit d'}\illegible{1} \unsure{les carrés|les images} \unsure{et antéimages|et anti-images} \illegible{1} \unsure{caractérisent|caractérisés} \unsure{un autre type} \unsure{supplémentaire|propre} \unsure{de structure|à ces structures}.}

Les foncteurs $u^*:\mathrm{HO}_W(I)\unsure{-|\to}\mathrm{HO}_W(I')$
% archive sheet number: 60
\scanpage{62}
\mspage{55}

sont ``homotopiquement exacts'' dans le sens de ces structures ``triangulées'' (\unsure{sans commentaires|aux commentaires|sans axiome de Verd.}), i.e. transforment suites exactes de susp. resp. de cosuspension en \unsure{itou|idem}.

Quels sont les axiomes des dérivateurs qui sont vérifiés \unsure{pour|sous} la seule hypothèse quillénisable ? $\mathrm{Der}$ (1) l'est sans hypothèse aucune sur la cat. de modèles $(M,W)$ :

\[ \deleted{\unsure{\mathrm{Der}}}\qquad\mathrm{HO}_{\deleted{\unsure{M,}}W}(\coprod I_\alpha)\simeq\prod_\alpha\mathrm{HO}_W(I_\alpha) \]

pour une famille \uline{finie} d'objets $I_\alpha$ dans $\mathrm{Diag}$. L'axiome de localisation \inserted{$\mathrm{Der}$ (2)} l'est sous l'hypothèse que $W$ soit \uline{fortement} \uline{saturé}, et quillénisable. On utilise le fait que \unsure{toute} flèche de $\mathrm{HO}_W(I)$ provient \deleted{\unsure{de}}, à isom. près, d'une flèche dans $\unsure{M(I)}$. (C'était là un des points qui \unsure{restait|restaient} en suspens dans $\overline{\underline{\mathrm{VI}}}$, dans la théorie \deleted{\unsure{g.}} des $\mathrm{HOT}_W$, $W$\inserted{$\subset\mathrm{Fl}(\mathrm{Cat})$} un loc. fond \ldots)

% slanted

\margin{Il est naturel ici de supposer que $W$ satisfait à \deleted{\unsure{tous}} les ``bons axiomes'', \unsure{sur les localiseurs fondamentaux|i.e. ceux du loc. fondamental} $W(\unsure{1},2,3,4,5,6,7)$.}

L'axiome $\mathrm{Der}$ 3, nouvelle mouture, implique un loc fondamental $\mathbb{W}\subset\mathrm{Fl}(\mathrm{Cat})$ (-- pas forcément $\mathbb{W}_\infty$) La forme la plus forte serait celle ci : Si

\[ u:I'\unsure{-|\to}I \]
% archive sheet number: 61
\scanpage{63}
\mspage{56}

est dans $\mathbb{W}$, on veut que le foncteur $u^*$ induise une \uline{équivalence}

\[ u^{*}_{\mathrm{l.c.}}:\mathrm{HO}^{\unsure{\mathrm{lc}|\ell\mathrm{c}}}_W(I)\xrightarrow{\approx}\mathrm{HO}^{\mathrm{lc}}_W(I'). \]

Mais je doute qu'on soit armé pour vérifier un axiome aussi fort, sans avoir \inserted{au moins} une théorie des opérations $u_!$ ou $u_*$.

On pourrait à \deleted{tout} hasard poser cet axiome \deleted{\unsure{pour}} en prenant $\mathbb{W}=\mathbb{W}_{\unsure{\alpha|a}}\overset{\text{déf}}{\unsure{\ne|=}}$ le plus petit loc. fond satisfaisant $W(1\text{ à }7)$ (le plus grand de ces localiseurs étant $\mathbb{W}_\infty$). Pour l'instant, la façon de vérifier cet axiome \deleted{\unsure{(qu}} (en termes de la donnée $(M,W)$) est pour moi un mystère $\pm$ total. (Mais \unsure{le|je} sais le vérifier pour $\mathrm{HOT}_W$, quand $W\subset\mathrm{Fl}(\mathrm{Cat})$ satisfait $W(1,2,3,4,5\ldots)$

Je vais m'en occuper, une fois que j'aurai une maîtrise raisonnable sur le formalisme de $u_*,u_!$ relatif \unsure{aux|au} $u:I'\to I$ dans $\mathrm{Diag}$. L'hypothèse quillénisable \inserted{seule} semble trop grossière pour assurer l'existence de $u_{\unsure{!|*}}$ ou de $u_*$ dans des cas intéressants.
% archive sheet number: 62
\scanpage{64}
\mspage{57}

6. Cat. de modèles quillénisables, et foncteurs $u_!$, $u_*$.

\uline{Définition} 1 : \textcircled{a} On dit que $(M,W)$ \deleted{est} \uline{uniformément quillénisable} à gauche, \inserted{sur $\mathrm{Diag}$} si on peut trouver une str. de Qu f $(C,\mathrm{TC},F,\mathrm{TF})$ \deleted{pour} $M$ qui \inserted{soit \uline{close}, et qui de plus} ait la propriété suivante : pour $I$ dans \deleted{$\mathrm{Diag}$} \inserted{$\mathrm{Cat}$}, \deleted{soient} \deleted{\unsure{a}} $C(I),\mathrm{TC}(I)\subset\mathrm{Fl}(M(I))$ définis par

\[
\begin{aligned}
(u:X\longrightarrow Y)\in C(I)\quad(\text{resp. }\mathrm{TC}(I))\quad&\text{ \unsure{sss|ssi} }\forall s\in\mathrm{Ob}\,I,\\
u_s:X_s\longrightarrow \unsure{Z|Y}_s&\text{ est }\in C(I)\quad(\text{resp. }\mathrm{TC}(I)).
\end{aligned}
\]

Soit de plus

\[ F(I)=\mathrm{TC}(I)_{\unsure{*|\alpha}},\qquad \mathrm{TF}(I)=C(I)_*, \]

La condition envisagée, c'est que pour \deleted{\unsure{tt}} $I$ \uline{dans $\mathrm{Diag}$}, $(\unsure{T|C}(I),\mathrm{TC}(I),F(I),\mathrm{TF}(I))$ est une structure de Qu f \uline{close} \uuline{exacte} \deleted{pour} $\deleted{\unsure{T|\mathrm{TF}}}W(I)$.$^{*}$/

% margin note written sideways

\margin{$^{*}$ $(C,\mathrm{TC},F,\mathrm{TF})$ est appelé alors une \uline{quillénisation uniforme}, à g. de $(M,W)$ sur $\mathrm{Diag}$.}

\textcircled{b} On dit que $(M,W)$ est quillénisable sur $\mathrm{Diag}$ uniformément et \inserted{\unsure{exactement|exactemt}} \unsure{gauche|à gauche}, si on \deleted{\unsure{pr}} a de plus $\mathrm{TC}\unsure{=|\ne}C\cap W$, i.e. si $(C,\mathrm{TC},F,\mathrm{TF})$ \deleted{est} exacte à gauche, donc définie par $C$ seule via

\[ \mathrm{TC}=C\cap W,\quad F=(\deleted{\unsure{\mathrm{TC}|\mathrm{TF}}}\unsure{C\cap W|\mathrm{TC}})_*,\quad \mathrm{TF}=C_*.\ /{}^{**} \]

% margin note written sideways

\margin{$^{**}$ On dit que $(C,\ldots)$, ou simplement $C$, définit une quillénisation uniforme et exacte à g. de $(M,W)$ sur $\mathrm{Diag}$.}

On a vu au \S~3 que si $M$ est \deleted{accessible} stable par petites $\underset{\unsure{-|\to}}{\lim}$ et accessible, et $W$ stable par petites $\varinjlim$ filtrantes, enfin si
% archive sheet number: 63
\scanpage{65}
\mspage{58}

$W$ admet une structure de \unsure{Qu.f.} \inserted{$(C,\mathrm{TC},F,\mathrm{TF})$} qui soit close et exacte à gauche -- alors les conditions de la définition 1 sont satisfaites. \deleted{\unsure{La p}}

D'autre part, notons que la prop. 1 page \unsure{45|55}, et ses corollaires, y inclus la prop 3 (p. 50) et son corollaire, sont \deleted{\unsure{satisfait}} valables dès lors que les quillénisations des $(M(I),W(I))$ envisagées sont définies ``uniformément à gauche'' par une quillénisation uniforme à gauche (pas \unsure{néct|nécess.} ``exacte à gauche'') de $(M,W)$ sur $\mathrm{Diag}$.

La quillénisation uniforme gauche nous servira à établir l'existence des images directes homotopiques $u_{\unsure{*}}$, pour certains $u$. La quillénisation uniforme droite, \unsure{dualement}, \unsure{quand|que} nous \unsure{considérons|considérerons} \unsure{principalt|principalement|présentement} l'\unsure{existce|existence} des $u_{\unsure{!|*}}$ \inserted{homotopiques}, pour certains $u:I'\to I$ dans $\mathrm{Diag}$.

Je vais résumer ce que j'ai pu établir jusqu'à présent dans ce sens, dans \unsure{la|le|les}
% archive sheet number: 64
